\originalchapter{传统重整化群、运行耦合与渐近自由}
\label{ch:traditional-rg}

Wilson 作用量把高能自由度的影响写进随截止变化的耦合.传统微扰计算通常不显式积分一个个壳层,而是在减除发散时引入尺度 $\mu$.同一个裸理论可以在不同 $\mu$ 下被分成“重整化作用量加反项”;如果分法改变而物理不变,重整化参数就必须随 $\mu$ 协同变化.本章从这个要求推导流方程,再用它处理大对数、QED 的屏蔽和 QCD 的渐近自由.

\originalsection{减除方案与重整化尺度}

消去极点只固定反项的发散部分.有限部分可以用在壳条件、指定动量点的顶角,或其他相容条件定义.这些选择称为重整化方案.在壳方案把参数尽量定义为极点质量和低动量实验耦合;最小减除(MS)只减 $1/\epsilon$ 极点;修正最小减除($\overline{\rm MS}$)还减去伴随的 $-\gamma_E+\log4\pi$.本书采用 $d=4-\epsilon$ 时,常见一圈因子是
\begin{equation}
 \mathcal D_\epsilon=\frac2\epsilon-\gamma_E+\log4\pi.
\end{equation}
若改用 $d=4-2\varepsilon$,同一个极点写作 $1/\varepsilon$.不能只搬用教材中的系数而不搬用其维数约定.在六维三次理论中,同样取 $d=6-\epsilon$,相应 Gamma 函数也在 $\epsilon/2$ 处展开.

同一裸场可以有 $\phi_B=Z_{\rm OS}^{1/2}\phi_{\rm OS}
=Z_{\overline{\rm MS}}^{1/2}\phi_{\overline{\rm MS}}$.两个 $Z$ 分别发散,但其比值在相容的紫外减除后有限;对于带电在壳场还要保留同一个红外调节器.不同方案的耦合也通过有限重定义联系.改变这些参数的定义不改变从完整振幅抽取的实验结果,尽管有限阶截断仍有高阶的方案差异.

为什么需要在壳方案之外的方案？无质量外腿的在壳减除可能被红外奇点阻碍;而在 $Q\gg m$ 的过程中,即使 $\alpha$ 很小,$\alpha\log(Q/m)$ 也可能不是小量.用低能耦合计算高能振幅会留下大对数,选 $\mu\sim Q$ 可以把显式对数缩小,将尺度差移入耦合运行.这里 $\mu$ 是计算和参数定义的尺度,不是新的物理阈值.

\originalsection{从裸量不变性提取 beta 函数}

定义重整化参数的 beta 函数时,必须注明保持裸量不变:
\begin{equation}
 \beta_g=\left.\mu\frac{\dd g}{\dd\mu}\right|_B.
 \label{eq:traditional-beta-definition}
\end{equation}
以一个在整数临界维数无量纲的耦合 $u$ 为例,在 MS 中可写
\begin{equation}
 u_B=\mu^{a\epsilon}
       \left[u+\frac{c_1(u)}\epsilon+
                    \frac{c_2(u)}{\epsilon^2}+\cdots\right].
 \label{eq:bare-pole-series}
\end{equation}
这里 $a$ 由耦合在延拓维数中的工程维数确定.四维四次标量耦合取 $a=1$,QED 的 $e$ 取 $a=1/2$,$\alpha$ 取 $a=1$.令
$\beta_\epsilon=-a\epsilon u+\beta(u)$,对 \eqref{eq:bare-pole-series} 求 $\mu$ 导数,有限阶的 $\epsilon^0$ 项给
\begin{equation}
 \beta(u)=a\left[u c_1'(u)-c_1(u)\right].
 \label{eq:beta-from-simple-pole}
\end{equation}
极点项还给更高阶 $c_n$ 的一致性条件.不能简单令 $\epsilon=0$ 后再求导:$-a\epsilon u$ 乘以 $1/\epsilon$ 正是产生有限 beta 函数的项.

若有多个耦合且工程维数为 $\delta_i+a_i\epsilon$,写
$g_{iB}=\mu^{\delta_i+a_i\epsilon}[u_i+C_i^{(1)}(u)/\epsilon+\cdots]$,相应有限项为
\begin{equation}
 \beta_i=-\delta_i u_i-a_iC_i^{(1)}
              +\sum_j a_j u_j\frac{\partial C_i^{(1)}}{\partial u_j},
 \label{eq:multi-beta-simple-pole}
\end{equation}
其余极点项必须相互抵消.$\delta_i$ 非零时,维数齐次性也进入这些条件;这不是对任意不满足局域量纲的极点系数都成立的形式规则.

\begin{example}
四维 $\lambda\phi^4/4!$ 理论的一圈关系为
$\lambda_B=\mu^\epsilon[\lambda+3\lambda^2/(16\pi^2\epsilon)+\cdots]$.求 beta 函数并与壳层流比较.
\end{example}
\begin{solution}
取 $a=1$,$c_1=3\lambda^2/(16\pi^2)$,由 \eqref{eq:beta-from-simple-pole} 得
\begin{equation}
 \beta_\epsilon(\lambda)=-\epsilon\lambda+
                      \frac{3\lambda^2}{16\pi^2}+\order(\lambda^3).
 \label{eq:phi4-traditional-beta}
\end{equation}
前章 $g=K_4\lambda=\lambda/(8\pi^2)$,所以
$\Lambda\dd g/\dd\Lambda=-\epsilon g+3g^2/2$ 在临界质量区成立.改变为 $t=\log(\Lambda_0/\Lambda)$ 后正好成为 $\dot g=\epsilon g-3g^2/2$.非零质量时,Wilson 壳层流保留阈值分母,而质量无关的 MS beta 函数不显式含这些分母;不能因此把两种作用量在任意质量、任意外动量下逐项等同.
\end{solution}

在方案重定义 $\widetilde u=f(u)$ 下,链式法则给
$\widetilde\beta(f(u))=f'(u)\beta(u)$,与 Wilson 耦合空间的向量规律相同.若单个边缘耦合在临界维数满足
$\beta=b_1u^2+b_2u^3+\cdots$,且 $f(u)=u+a_2u^2+a_3u^3+\cdots$,展开并用 $u=\widetilde u-a_2\widetilde u^2+\cdots$ 代回,得
\begin{equation}
 \widetilde\beta=b_1\widetilde u^2+b_2\widetilde u^3+\order(\widetilde u^4).
\end{equation}
前两个系数在这种限定下不变;高阶系数一般会变.“前两个系数普适”不是对任意多耦合体系、任意质量相关方案的无条件陈述.在固定点,可逆坐标变换使稳定矩阵相似,因而本征值不变;这解释了临界指数为何比固定点的坐标值更直接具有物理意义.

\originalsection{六维三次理论:质量运行与维数嬗变}

原课用六维实三次标量理论演示 beta 函数.它也能说明负 beta 函数的效果,而不需要先计算非 Abel 图.这里令 $d=6-\epsilon$,考虑 Euclidean 微扰作用量
\begin{equation}
 S=\int\dd^dx\left[\frac12(\partial\phi)^2+\frac12m^2\phi^2
                    +\frac{\mu^{\epsilon/2}g}{3!}\phi^3\right],
 \qquad a=\frac{g^2}{(4\pi)^3}.
\end{equation}
实 $\phi^3$ 势没有下界,因此这是围绕局部真空的形式微扰示例,不给出稳定非微扰真空的构造;它不能作为与 QCD 同等的完整物理模型.

把反项系数定义为动能 $1+A$、质量项 $m^2(1+B)$、三点项 $g(1+C)$.一圈二点积分含
\begin{equation}
 I_2(p)=\frac{\mu^\epsilon\Gamma(2-d/2)}{(4\pi)^{d/2}}
       \int_0^1\dd x\,[m^2+x(1-x)p^2]^{d/2-2}.
\end{equation}
$\Gamma(-1+\epsilon/2)=-2/\epsilon+\cdots$,且
$\int_0^1x(1-x)\dd x=1/6$.
Euclidean 一粒子不可约二点核的环修正为 $-g^2I_2/2$,故极点为 $a(m^2+p^2/6)/\epsilon$,需要 $A=-a/(6\epsilon)$、$B=-a/\epsilon$.三点三角积分有
$\Gamma(\epsilon/2)\int_{x+y+z=1}\dd x\dd y$,单纯形面积 $1/2$,给顶角相对极点 $a/\epsilon$,因而 $C=-a/\epsilon$.这些系数的减除约定是 MS;有限部分由方案另定.

裸量关系必须包含场归一化的逆因子:
\begin{equation}
 \phi_B=(1+A)^{1/2}\phi,\quad
 m_B^2=m^2\frac{1+B}{1+A},\quad
 g_B=\mu^{\epsilon/2}g\frac{1+C}{(1+A)^{3/2}}.
 \label{eq:phi3-bare-relations}
\end{equation}
因此
\begin{equation}
 a_B=\mu^\epsilon\left[a-\frac{3a^2}{2\epsilon}+\cdots\right],
 \qquad m_B^2=\mu^2r\left[1-\frac{5a}{6\epsilon}+\cdots\right],
 \qquad r=\frac{m^2}{\mu^2}.
\end{equation}
由简单极点公式得
\begin{equation}
 \beta_a=-\frac32a^2+\order(a^3),\qquad
 \beta_r=-2r-\frac56ar+\order(a^2r),\qquad
 \gamma_\phi\equiv\frac12\mu\frac{\dd\log Z_\phi}{\dd\mu}
       =\frac a{12}+\order(a^2),
 \label{eq:phi3-six-d-beta}
\end{equation}
这里已取 $\epsilon\to0$.第一条为负,所以耦合向紫外减小、向红外增大.第二条首先反映相关质量的工程缩放,而不是另一个边缘耦合.

设 $a(\mu_0)=a_0$,积分第一条得到
\begin{equation}
 \frac1{a(\mu)}=\frac1{a_0}+\frac32\log\frac\mu{\mu_0},\qquad
 a(\mu)=\frac{a_0}{1+\tfrac32a_0\log(\mu/\mu_0)}.
 \label{eq:phi3-running-alpha}
\end{equation}
同时有 $\dd\log m^2/\dd\log\mu=-5a/6$,与 $\dd a/\dd\log\mu=-3a^2/2$ 相除可积,得
\begin{equation}
 m^2(\mu)=m^2(\mu_0)\left[\frac{a(\mu)}{a_0}\right]^{5/9},\qquad
 r(\mu)=r(\mu_0)(\mu_0/\mu)^2
              \left[\frac{a(\mu)}{a_0}\right]^{5/9}.
\end{equation}
质量的工程增长与对数修正在此明确分开.

定义
\begin{equation}
 \Lambda=\mu_0\exp\left[-\frac2{3a_0}\right],\qquad
 a(\mu)=\frac2{3\log(\mu/\Lambda)}.
 \label{eq:phi3-dimensional-transmutation}
\end{equation}
对 $\mu_0$ 求导并使用 beta 函数,可验证 $\Lambda$ 与参考点无关.原来无量纲耦合加参考尺度的输入,可换成一个有量纲尺度,这称为维数嬗变.它是微扰流所给的尺度,不是由一圈方程证明的束缚态质量.$\mu$ 接近 $\Lambda$ 时 $a$ 变大,一圈近似在形式极点之前已经失效.

展开 \eqref{eq:phi3-running-alpha} 给
\begin{equation}
 a(\mu)=a_0\sum_{n=0}^\infty
       \left[-\frac32a_0\log\frac\mu{\mu_0}\right]^n
\end{equation}
在几何级数收敛区成立.高能时原来的固定阶展开会被大对数破坏,但 beta 方程只要求沿途每个无穷小尺度步中的耦合小;它把主对数系列组织为有理函数.重求和不等于得到了全部高圈系数,更不允许越过强耦合区域继续声称可控.

\originalsection{Callan--Symanzik 方程与特征线}

设 $G_n(x_1,\ldots,x_n;g,m^2,\mu)$ 是重整化场的连通关联函数,取点互不重合以避免额外复合算符接触项.裸函数满足
\begin{equation}
 G_{nB}=Z_\phi^{n/2}G_n.
\end{equation}
保持裸量不变时,左侧不依赖 $\mu$,因此
\begin{equation}
 \left[\mu\frac\partial{\partial\mu}
       +\beta_g\frac\partial{\partial g}
       +\beta_{m^2}\frac\partial{\partial m^2}
       +n\gamma_\phi\right]G_n=0.
 \label{eq:callan-symanzik-connected}
\end{equation}
这是一阶偏微分方程;$\mu\partial_\mu$ 只作用于显式尺度,beta 项补上参数随尺度的变化.若改用 $r=m^2/\mu^2$,则 $\beta_r=\beta_{m^2}/\mu^2-2r$,并相应改变显式导数的保持量.对一粒子不可约 $n$ 点顶角,场因子倒过来,方程中的 $n\gamma_\phi$ 变为 $-n\gamma_\phi$,不能与连通函数混用.

令 $s=\log(\mu'/\mu)$,沿特征线解
$\dd g/\dd s=\beta_g$、$\dd m^2/\dd s=\beta_{m^2}$,则偏微分方程化为
$\dd G_n/\dd s=-n\gamma_\phi(g(s))G_n$.因此
\begin{equation}
 G_n(x;g(\mu'),m^2(\mu'),\mu')
 =\exp\left[-n\int_\mu^{\mu'}\frac{\dd\nu}\nu
                   \gamma_\phi(g(\nu))\right]
 G_n(x;g(\mu),m^2(\mu),\mu).
 \label{eq:cs-characteristic-solution}
\end{equation}
若只有一个无质量耦合,积分也可写成 $\int\gamma_\phi(g)\dd g/\beta(g)$.这就是场归一化的运行.关联函数本身不必保持数值不变,因为场的定义变了;由适当外腿留数、源与顶角组成的可观测量才不依赖 $\mu$.

\begin{example}
推导无质量标量二点函数的高动量缩放,并在固定点读出场的维数.
\end{example}
\begin{solution}
去掉动量守恒 delta 函数的二点函数有工程维数 $-2$,所以
$G_2(p;g,\mu)=p^{-2}f(p/\mu,g)$.量纲缩放给
$G_2(sp;g(s\mu),s\mu)=s^{-2}G_2(p;g(s\mu),\mu)$.
再用 \eqref{eq:cs-characteristic-solution} 将左侧的尺度从 $s\mu$ 改回 $\mu$,得
\begin{equation}
 G_2(sp;g(\mu),\mu)=s^{-2}
 \exp\left[2\int_\mu^{s\mu}\frac{\dd\nu}\nu\gamma_\phi(g(\nu))\right]
 G_2(p;g(s\mu),\mu).
 \label{eq:cs-two-point-momentum-scaling}
\end{equation}
在固定点 $g=g^*$,$\gamma_\phi=\gamma_\phi^*$ 为常数,故动量幂次为
$G_2(sp)=s^{-2+2\gamma_\phi^*}G_2(p)$.
Fourier 变换对应的坐标空间维数为
\begin{equation}
 \Delta_\phi=\frac{d-2}{2}+\gamma_\phi^*,\qquad
 G_2(x)\propto\frac1{|x|^{2\Delta_\phi}},\qquad
 \eta=2\gamma_\phi^*.
 \label{eq:cs-fixed-point-dimension}
\end{equation}
也可从坐标空间
$G_2(x)=\mu^{d-2}f(\mu x)$ 直接求出
$[y\partial_y+d-2+2\gamma_\phi^*]f=0$.
固定 $\mu$ 下比例系数依赖场归一化;缩放幂次是核心结论.复合算符有自己的重整化和混合矩阵,不能全部用 $n\Delta_\phi$ 猜出其维数.
\end{solution}

\originalsection{QED beta 函数与有效电荷}

Ward 恒等式给 $e_B=\mu^{\epsilon/2}Z_3^{-1/2}e$.前章一圈光子反项给
$Z_3=1-2\alpha/(3\pi\epsilon)+\cdots$,因此
\begin{equation}
 \alpha_B=\mu^\epsilon\left[\alpha+
              \frac{2\alpha^2}{3\pi\epsilon}+\cdots\right].
\end{equation}
简单极点公式立即给出一个单位电荷 Dirac 物种的结果
\begin{equation}
 \beta_\alpha=\frac{2\alpha^2}{3\pi}+\order(\alpha^3),\qquad
 \beta_e=\frac{e^3}{12\pi^2}+\order(e^5).
 \label{eq:qed-beta-one-loop}
\end{equation}
若有活跃物种 $Q_i e$,右侧乘 $\sum_iQ_i^2$,并计入内部简并度.这里的“活跃”需要结合能量相对质量讨论,不能在任意低能都数入所有重粒子.

取 $\alpha(\mu_0)=\alpha_0$,积分得到
\begin{equation}
 \frac1{\alpha(\mu)}=\frac1{\alpha_0}-\frac2{3\pi}\log\frac\mu{\mu_0},\qquad
 \alpha(\mu)=\frac{\alpha_0}{1-\tfrac{2\alpha_0}{3\pi}\log(\mu/\mu_0)}.
 \label{eq:qed-running-alpha}
\end{equation}
向紫外耦合增加,向红外耦合减小.若无质量电子一直适用,红外趋向自由固定点,电荷是边缘无关的.形式上的紫外 Landau 尺度为
\begin{equation}
 \Lambda_L=\mu_0\exp\left[\frac{3\pi}{2\alpha_0}\right].
\end{equation}
这标记一圈外推的失效尺度,不是可用微扰论描述的真实无限耦合状态;实际 QED 在更高能会嵌入其他物理,也不能据此一圈式证明或否定其严格连续极限.

有质量电子时,实验有效电荷还依赖动量.采用前章 $1-\Pi$ 的分母约定,$\overline{\rm MS}$ 有
\begin{equation}
 \Pi_{\overline{\rm MS}}(-Q^2;\mu)
   =\frac{2\alpha(\mu)}\pi\int_0^1\dd x\,x(1-x)
             \log\frac{m^2+x(1-x)Q^2}{\mu^2}+\order(\alpha^2).
\end{equation}
守恒电流交换定义
\begin{equation}
 \widehat\alpha(Q)=\frac{\alpha(\mu)}
                 {1-\Pi_{\overline{\rm MS}}(-Q^2;\mu)}.
 \label{eq:qed-effective-ms-alpha}
\end{equation}
其 $\mu$ 导数在所算阶数消失:$\partial_{\log\mu}\Pi=-2\alpha/(3\pi)$,正好抵消 $\beta_\alpha=2\alpha^2/(3\pi)$ 对分子的贡献.若忽略未计算高圈项,保留分母形式只是对泡链对数的组织,不代表已经计算所有 $\alpha^n$ 项.

在 $Q=0$,取本阶 $m$ 为电子极点质量,
\begin{equation}
 \Pi_{\overline{\rm MS}}(0;\mu)=\frac{2\alpha(\mu)}{3\pi}\log\frac m\mu,
 \qquad
 \widehat\alpha(0)=\frac{\alpha(\mu)}
        {1-\tfrac{2\alpha(\mu)}{3\pi}\log(m/\mu)}.
\end{equation}
与 \eqref{eq:qed-running-alpha} 比较,在此一圈精度下
$\widehat\alpha(0)=\alpha(\mu=m)$.这个等式依赖所用减除约定,高阶匹配会修正它.
对 $Q\gg m$,$\Pi\sim\alpha\log(Q^2/\mu^2)/(3\pi)$,选 $\mu\sim Q$ 后 $\widehat\alpha(Q)$ 与运行耦合相近;两者还可有方案依赖的有限差.

对 $Q\ll m$,物理真空极化按 $Q^2/m^2$ 消失,$\widehat\alpha(Q)$ 趋于固定的低能电荷.但在质量无关的 MS 方程中,$\alpha(\mu)$ 若被机械跑到 $\mu\ll m$,仍按无质量物种的 beta 函数变化.这是参数定义的性质,不是实验电荷真的趋于零.正确的低能做法是在 $\mu\sim m$ 匹配:积分掉电子,得到 Maxwell 理论及 $F\Box F/m^2$、$F^4/m^4$ 等高维项;纯 Maxwell 理论没有带电环,最低维电荷不再运行.高维项在 $Q/m$ 中受抑制.这就是阈值退耦,把 Wilson 低能自由度图景与传统方案联系起来.

\originalsection{QCD 的一圈系数与反屏蔽}

对 $SU(N_c)$,取 Hermitian 生成元满足
$\tr(t^at^b)=T_F\delta^{ab}$,基本表示 $T_F=1/2$,伴随 Casimir $C_A=N_c$.采用与 QED 一致的 $D_\mu=\partial_\mu+\ii gA_\mu^at^a$;于是
\begin{equation}
 F_{\mu\nu}=\partial_\mu A_\nu-\partial_\nu A_\mu+\ii g[A_\mu,A_\nu],\qquad
 \Lag=-\frac14F_{\mu\nu}^aF^{a\mu\nu}
      +\sum_{j=1}^{N_f}\bar\psi_j(\ii\slashed D-m_j)\psi_j.
\end{equation}
换用 $D=\partial-\ii gA$ 的约定只改变相应顶角和场强的符号约定,不改变以下 $g^2$ 的流.

非 Abel 场本身带颜色,因此除了夸克环,还有规范场与鬼场环.一圈系数不能从 QED 的电荷数目直接类比得到.为了显示这几个贡献怎样组成,可把耦合吸收到背景连接中,使 Euclidean 经典作用为
$S_B=\tfrac1{4g^2}\int(F_B^a)^2$.此小节中的下标 $B$ 表示背景(background),不表示前文的裸量;$F_B^a$ 是已经吸收耦合 $g$ 的背景曲率分量.
取反 Hermitian 连接 $\mathcal B_\mu=\ii gA_\mu^at^a$,$\mathcal F_{\mu\nu}=[D_\mu,D_\nu]$.在背景 Feynman 规范中,二次涨落的向量和鬼场算符分别为
\begin{equation}
 (\Delta_1)_{\mu\nu}=-D^2\delta_{\mu\nu}-2\,\operatorname{ad}\mathcal F_{\mu\nu},
 \qquad \Delta_0=-D^2.
\end{equation}
其 Gaussian 积分给一圈作用
$\Gamma_1=\tfrac12\Tr\log\Delta_1-\Tr\log\Delta_0$.
每个无质量 Dirac 物种还给
$-\tfrac12\Tr\log\Delta_{1/2}$,其中
$\Delta_{1/2}=-D^2-\tfrac12\gamma^{\mu\nu}\mathcal F_{\mu\nu}$,
$\gamma^{\mu\nu}=[\gamma^\mu,\gamma^\nu]/2$ 是 Euclidean 伽马矩阵.

下面调用一个解析工具,而不在本课程内证明其一般形式:对平直、无边界背景中的 Laplace 型算符
$\Delta=-(D^2+E)$,热核小 $s$ 展开中决定四维对数发散的积分系数是
\begin{equation}
 A_4=\int\dd^4x\,\tr\left(\frac12E^2+
                  \frac1{12}\Omega_{\mu\nu}\Omega_{\mu\nu}\right),
 \qquad \Omega_{\mu\nu}=[D_\mu,D_\nu].
 \label{eq:qcd-heat-kernel-tool}
\end{equation}
总导数已在无边界条件下积分掉.
其条件与归一化见 Vassilevich 的热核综述公式 (4.28);该文的平直 Yang--Mills 代入例为 (4.30)--(4.34).\footnote{D.~V.~Vassilevich, \emph{Heat kernel expansion: user's manual}, arXiv:hep-th/0306138v3,\url{https://arxiv.org/abs/hep-th/0306138},印刷页 40--41.这里仅引用平直空间的一圈解析系数;没有调用或重述该文其他几何结果.}

现在的群与自旋迹可在本地核算.把 $X=\tr_R(\mathcal F_{\mu\nu}\mathcal F_{\mu\nu})$ 作为中间记号.对向量算符,连接曲率作用在四个 Lorentz 分量上,$\tr\Omega^2=4X$;$E_{\mu\nu}=2\mathcal F_{\mu\nu}$,故
$\tr E^2=4\tr_R(\mathcal F_{\mu\nu}\mathcal F_{\nu\mu})=-4X$.
于是向量的 $A_4$ 密度为 $-5X/3$,鬼场为 $X/12$.按行列式的 $1/2$ 与 $-1$ 权重相加,得到
$-11\tr_{\rm adj}\mathcal F^2/12$.
对于 Dirac 算符,$\tr\Omega^2=4X$;用
\begin{equation}
 \tr_{\rm spin}(\gamma^{\mu\nu}\gamma^{\rho\sigma})
 =4(\delta^{\mu\sigma}\delta^{\nu\rho}
             -\delta^{\mu\rho}\delta^{\nu\sigma})
\end{equation}
可得 $\tr E^2=-2X$,所以 $A_4$ 密度为 $-2X/3$;乘费米子权重 $-1/2$ 后为 $X/3$.
反 Hermitian 生成元的迹满足
$\tr_R\mathcal F^2=-T_R(F_B^a)^2$,伴随表示取 $T_R=C_A$.总加权系数因此为
\begin{equation}
 \left(\frac{11C_A}{12}-\frac{N_fT_F}{3}\right)(F_B^a)^2
       =\frac{b_0}{4}(F_B^a)^2,\qquad
 b_0=\frac{11}{3}C_A-\frac43T_FN_f.
 \label{eq:qcd-b0-trace-computation}
\end{equation}
这段计算把 $11/3$ 的规范场加鬼场贡献与 $-4T_FN_f/3$ 的费米子屏蔽明确分开;热核工具本身的解析证明是已说明的外部输入.

用 $\Tr\log\Delta=-\int\dd s\,\Tr\e^{-s\Delta}/s$,四维热核前因子是 $(4\pi s)^{-2}$;在 $1/\Lambda^2<s<1/\mu^2$ 上的对数为 $2\log(\Lambda/\mu)$.故
\begin{equation}
 \Delta S=-\frac{b_0}{32\pi^2}\log\frac\Lambda\mu
                \int\dd^4x(F_B^a)^2,
 \qquad
 \frac1{g^2(\mu)}=\frac1{g^2(\Lambda)}
                   -\frac{b_0}{8\pi^2}\log\frac\Lambda\mu.
\end{equation}
对 $\mu$ 求导,得到
\begin{equation}
 \beta_g=-\frac{g^3}{16\pi^2}
         \left(\frac{11}{3}N_c-\frac23N_f\right)+\order(g^5),
 \qquad
 \beta_{\alpha_s}=-\frac{b_0}{2\pi}\alpha_s^2+\order(\alpha_s^3),
 \quad\alpha_s=\frac{g^2}{4\pi}.
 \label{eq:qcd-beta-one-loop}
\end{equation}
对 $N_f<11N_c/2$,$b_0>0$,耦合随紫外尺度上升而减小,称为渐近自由.规范场反屏蔽超过了夸克的屏蔽,是与 QED 相反的关键机制.

\originalsection{运行、阈值与强耦合边界}

在活跃夸克数固定的能量区间,一圈积分给
\begin{equation}
 \frac1{\alpha_s(\mu)}=\frac1{\alpha_s(\mu_0)}
                  +\frac{b_0}{2\pi}\log\frac\mu{\mu_0},\qquad
 \alpha_s(\mu)=\frac{2\pi}{b_0\log(\mu/\Lambda_{\rm QCD})},
 \label{eq:qcd-one-loop-running}
\end{equation}
其中
\begin{equation}
 \Lambda_{\rm QCD}=\mu_0
   \exp\left[-\frac{2\pi}{b_0\alpha_s(\mu_0)}\right].
\end{equation}
与六维三次例子一样,$\Lambda_{\rm QCD}$ 在本阶与参考点无关,但会随重整化方案和高阶定义改变.遇到重夸克阈值,应在质量附近匹配不同 $N_f$ 的有效理论,而不是从最高能到最低能始终固定 $N_f=6$.

\begin{center}
\begin{tikzpicture}[>=Stealth,scale=1]
 \draw[->](0,0)--(5.7,0) node[right] {$\log\mu$};
 \draw[->](0,0)--(0,3.1) node[above] {$1/\alpha$};
 \draw[thick,structurecolor](.5,.6)--(4.9,2.6) node[right] {QCD};
 \draw[thick](.5,2.6)--(4.9,.6) node[right] {QED};
 \node[align=center] at (2.7,-.7){同一活跃粒子区间内的一圈示意\\截距和尺度范围不作数值比较};
\end{tikzpicture}
\end{center}

在 $\mu\gg\Lambda_{\rm QCD}$ 时,如果过程也没有额外的大红外或运动学对数,$\alpha_s(\mu)$ 小,可以用短距离微扰.若尺度下降到 $\Lambda_{\rm QCD}$ 附近,耦合增大,原先的夸克、胶子弱耦合展开失效.式中形式极点不能被解释为真实散射振幅在同一点无穷大,也不是禁闭定理.

QCD 的低能现象包括颜色禁闭和近似手征对称性的自发破缺.对足够轻的夸克,在忽略质量与相关异常限制的条件下,经典作用具有左右手味变换;真空可有 $\langle\bar q q\rangle\ne0$,将左右手联系起来,并产生轻的赝 Goldstone 模式.非零质量使相应对称性只是近似的.这些陈述说明新的红外自由度为什么可能需要强子描述;它们不能从一圈 beta 函数单独推导.严格禁闭、手征凝聚值和谱,需要非微扰方法,超出本讲义的微扰证明范围.

从自由场到这里,重整化群补上了一条连续的关系:紫外子图要求局域反项,局域作用量的参数随分辨率改变,而这种改变既能把大对数重求和,也能指出弱耦合描述何时失效.可计算的内容是给定方案、粒子阈值和微扰精度下的流与振幅;新的强耦合物理必须另行求解,不能被一个运行公式替代.
